\section{Introduction}

Mixture-of-Experts (MoE) models~\cite{shazeer2017outrageously} scale neural-network capacity by routing each token to a small subset of expert subnetworks, keeping per-token compute bounded while the full parameter set must stay accessible for unconstrained routing. Modern MoE language models such as Mixtral~\cite{jiang2024mixtral}, DeepSeek-V2~\cite{deepseekv2}, and Qwen3~\cite{qwen3technicalreport} contain tens to hundreds of billions of parameters, yet activate only a few billion per token. Deploying these models on a single high-end GPU exposes a fundamental hardware constraint: device high-bandwidth memory (HBM), typically 48--80\,GB, must simultaneously hold expert weights, non-expert parameters, activations, the KV cache, and runtime buffers. Qwen3-Next-80B, for instance, requires roughly 160\,GB in FP16, far exceeding the HBM capacity of current accelerators. The deployment bottleneck is therefore not peak arithmetic throughput but \emph{on-device memory management}: deciding which expert weights remain resident, at what precision, and how residency adapts as the workload changes.

Two families of techniques address this memory-capacity wall.
The first, expert offloading and prefetching~\cite{li2023adaptive,yi2023edgemoe,xue2024moe,eliseev2023fast,hwang2024pre,zhong2024adapmoe,tang2024hobbit,song2024promoe,shen2025expertflow}, treats GPU HBM as a managed cache and moves experts across the PCIe memory hierarchy on demand.
The second, post-training quantization (PTQ)~\cite{lin2024awq,frantar2022gptq,kim2023mixturequantizedexpertsmoqe,duanmu2025mxmoe}, compresses expert weights to low bit-widths, reducing the per-expert device-memory footprint so the full pool fits on-chip.

Offloading works well when each iteration touches a small, stable expert subset. In practice, MoE activation becomes substantially denser during prefill and at larger batch sizes (\autoref{tab:expert_activation_ratio}), expanding the per-iteration working set beyond what can be staged within the PCIe compute-transfer overlap window. Once that window is exceeded, data movement stalls the GPU compute pipeline and amplifies worst-case latency (\autoref{fig:waiting-vs-prompt}).\looseness=-1

PTQ avoids cross-tier data movement entirely and, with no ongoing migration, achieves the lowest raw per-step latency. However, most PTQ pipelines fix a uniform or per-expert precision map offline. Expert utilization is heavy-tailed and task-dependent: the identity of frequently activated experts shifts across text, math, and code workloads (\autoref{fig:expert-active}). A static precision map therefore wastes scarce HBM capacity on experts that carry little traffic in the current task while over-compressing those that become critical. This paper does not claim to beat static PTQ on raw latency; it targets a better \emph{accuracy--throughput--memory} tradeoff under the same fixed device-memory budget.

We propose \systemname, which manages expert precision as a dynamically allocated on-device resource for single-GPU MoE inference under a declared HBM envelope. During inference, the system observes router outputs, estimates which experts carry disproportionate traffic over a stable time horizon, and allocates a limited high-precision budget to those experts. The remaining experts stay at lower precision. Preallocated pools and admission control keep all runtime-owned expert and transition memory within the envelope; the end-to-end cap holds when the deployment's KV-cache and activation reservations are respected.

The design separates policy from mechanism. A scheduler periodically selects budget-feasible high-precision sets per layer. A transition pipeline carries out promotions and demotions on a dedicated CUDA stream without blocking the forward path. Pool-based memory management with admission control prevents budget violations. Section~\ref{sec:design} describes each component.

The main contributions of this work are:
\begin{itemize}[leftmargin=1.5em, itemsep=1pt, topsep=2pt]
  \item We formulate single-GPU MoE inference under a hard device-memory constraint as an on-device precision residency management problem, targeting the accuracy--throughput--memory tradeoff rather than minimum raw latency alone.
  \item We design an asynchronous runtime execution model with versioned expert handles. A new handle is published only after its copy completes, and read leases plus per-compute-stream last-use events protect the old block until every dispatch and kernel that can reference it has drained.
  \item We introduce deterministic, pool-based device-memory management with admission control. Exact expert-sized blocks eliminate variable-size external fragmentation, per-tier capacities isolate resident slots, and a budget tracker prevents transient over-commitment during concurrent migrations.
  \item We evaluate \systemname\ on Qwen3-MoE-30B, Qwen3-Next-80B, and Phi-3.5-MoE across six benchmarks. \systemname\ recovers up to 79\% of the INT2-to-INT4 accuracy gap on Qwen3-Next-80B, from 73.58\% to 77.15\%. Performance experiments use uniform PTQ on all models and the pinned official MoE-Infinity runtime where its public implementation supports the evaluated model; ablations isolate each mechanism, and raw peak-memory traces test the 48\,GB deployment envelope.
\end{itemize}
